% Figure: how a reversible engine defines the thermodynamic temperature scale. A
% cascade of Carnot engines run between a sequence of reservoirs, each rejecting
% the heat the next absorbs. Because a reversible engine's heat ratio depends only
% on the reservoir temperatures, Q_k/Q_ref = T_k/T_ref, the heat exchanged is a
% linear measure of temperature independent of any working substance. The plot
% shows the heat ratio rising as a straight line through the origin against the
% ratio of thermodynamic temperatures, the defining property of the Kelvin scale.
\begin{figure}[htbp]
\centering
\begin{tikzpicture}
\begin{axis}[
  width=11cm, height=7cm,
  xlabel={reservoir temperature ratio $T/T_{\mathrm{ref}}$},
  ylabel={reversible heat ratio $Q/Q_{\mathrm{ref}}$},
  xmin=0, xmax=3, ymin=0, ymax=3,
  axis lines=left,
  tick label style={font=\scriptsize},
  label style={font=\small},
  ymajorgrids=true, xmajorgrids=true,
  grid style={enggray!22},
  legend style={font=\scriptsize, at={(0.03,0.97)}, anchor=north west, draw=enggray!40},
]
% The defining relation Q/Q_ref = T/T_ref: a straight line of slope one.
\addplot[engblue, very thick, domain=0:3, samples=2] {x};
\addlegendentry{$Q/Q_{\mathrm{ref}} = T/T_{\mathrm{ref}}$ (reversible)}
% cascade nodes: three reservoirs at ratios 0.5, 1.0, 2.0 with their heat ratios
\addplot[mark=*, only marks, engred, mark size=2.4pt]
  coordinates {(0.5,0.5) (1.0,1.0) (2.0,2.0)};
\node[font=\scriptsize, anchor=north west, engred] at (axis cs:2.0,2.0) {$T_3$};
\node[font=\scriptsize, anchor=north west, engred] at (axis cs:1.0,1.0) {$T_{\mathrm{ref}}$};
\node[font=\scriptsize, anchor=north west, engred] at (axis cs:0.5,0.5) {$T_1$};
% a hypothetical irreversible engine would fall below the line (rejects more heat)
\addplot[engorange, thick, dashed, domain=0:3, samples=2] {0.72*x};
\addlegendentry{irreversible engine (rejects more heat)}
\end{axis}
\end{tikzpicture}
\caption[The thermodynamic (Kelvin) temperature scale]{The reversible-engine
definition of temperature. A Carnot engine exchanging heat $Q$ with a reservoir at
$T$ and $Q_{\mathrm{ref}}$ with a reference reservoir at $T_{\mathrm{ref}}$ obeys
$Q/Q_{\mathrm{ref}} = T/T_{\mathrm{ref}}$, a ratio that depends on the two
temperatures alone and not on the working substance. Measuring the reversible heat
ratio therefore measures the temperature ratio directly, which is how the Kelvin
scale is fixed without reference to any particular thermometer fluid. The straight
line through the origin is that defining proportionality; the marked points are a
three-reservoir cascade in which each engine rejects the heat the next absorbs,
carrying the scale from one fixed point to another. A real irreversible engine
rejects more heat than the reversible relation allows and so falls below the line,
the same shortfall the Clausius inequality measures. Absolute zero is the
temperature at which a reversible engine rejects no heat at all, the origin where
the line meets the axis.}
\label{fig:vol04ch04:kelvin-scale}
\end{figure}
